\documentclass[11pt]{article}

%==========================================================
%  Packages
%==========================================================
\usepackage[utf8]{inputenc}
\usepackage[T1]{fontenc}
\usepackage[margin=1in]{geometry}
\usepackage{amsmath,amssymb,amsthm}
\usepackage{mathtools}
\usepackage{booktabs}
\usepackage{array}
\usepackage{enumitem}
\usepackage{xcolor}
\usepackage{tcolorbox}
\usepackage[colorlinks=true,linkcolor=blue!50!black,urlcolor=blue!50!black]{hyperref}

%==========================================================
%  Theorem-like environments
%==========================================================
\theoremstyle{definition}
\newtheorem{definition}{Definition}[section]
\newtheorem{assumption}[definition]{Assumption}
\newtheorem{problem}[definition]{Problem}

\theoremstyle{plain}
\newtheorem{theorem}[definition]{Theorem}
\newtheorem{lemma}[definition]{Lemma}
\newtheorem{proposition}[definition]{Proposition}

\theoremstyle{remark}
\newtheorem{remark}[definition]{Remark}

%==========================================================
%  Notation macros
%==========================================================
\newcommand{\Z}{\mathbb{Z}}
\newcommand{\R}{\mathbb{R}}
\newcommand{\Zq}{\Z_q}
\newcommand{\Zp}{\Z_p^*}
\newcommand{\vA}{\mathbf{A}}
\newcommand{\vB}{\mathbf{B}}
\newcommand{\vs}{\mathbf{s}}
\newcommand{\ve}{\mathbf{e}}
\newcommand{\vb}{\mathbf{b}}
\newcommand{\vx}{\mathbf{x}}
\newcommand{\vz}{\mathbf{z}}
\newcommand{\vzero}{\mathbf{0}}
\newcommand{\lat}{\mathcal{L}}
\newcommand{\dlog}{\mathrm{DLOG}}
\newcommand{\cdh}{\mathrm{CDH}}
\newcommand{\ddh}{\mathrm{DDH}}
\newcommand{\lwe}{\mathrm{LWE}}
\newcommand{\sis}{\mathrm{SIS}}
\newcommand{\svp}{\mathrm{SVP}}
\newcommand{\sivp}{\mathrm{SIVP}}
\newcommand{\gapsvp}{\mathrm{GapSVP}}
\newcommand{\bdd}{\mathrm{BDD}}
\newcommand{\norm}[1]{\lVert #1 \rVert}
\newcommand{\abs}[1]{\lvert #1 \rvert}
\newcommand{\lamone}{\lambda_1}
\newcommand{\reduces}{\;\leq\;}
\newcommand{\poly}{\mathrm{poly}}

% Boxed "key idea" environment
\newtcolorbox{keybox}[1][]{
  colback=blue!4!white, colframe=blue!50!black,
  fonttitle=\bfseries, title={#1},
  boxrule=0.5pt, arc=2pt, left=6pt,right=6pt,top=4pt,bottom=4pt}

%==========================================================
%  Title
%==========================================================
\title{\bfseries Discrete-Log vs.\ Lattice Assumptions\\[2pt]
\large A Comparison of Their Reduction Structures}
\author{Study Notes}
\date{\today}

\begin{document}
\maketitle

\begin{abstract}
\noindent
These notes compare two families of cryptographic hardness assumptions:
the \emph{discrete-logarithm / Diffie--Hellman} family (DLOG, CDH, DDH, and their
number-theoretic cousins) and the \emph{lattice} family (SVP/SIVP/GapSVP in the worst case,
and SIS/LWE in the average case). The focus is not on the problems themselves but on the
\emph{reductions between them} --- across dimensions, across parameters, and between
search and decision versions --- because it is precisely the reduction structure that
makes the lattice family robust and the DLOG family fragile. Four themes recur:
(i) lattice hardness \emph{propagates across dimensions} while DLOG hardness does not
propagate across moduli; (ii) both families enjoy \emph{random self-reducibility},
but for DLOG this is only a within-group loop that gives no guidance on choosing the group;
(iii) the average-case nature of DLOG rests on the \emph{Prime Number Theorem};
and (iv) \emph{search and decision are provably equivalent for LWE}, whereas in the
Diffie--Hellman world the decision problem is a separate, strictly stronger, and sometimes
\emph{false} assumption (pairing groups).
\end{abstract}

\tableofcontents
\bigskip
\hrule

%==========================================================
\section{Preliminaries}
%==========================================================

We fix notation and recall the problems that appear throughout.

\subsection{Discrete-logarithm family}

\begin{problem}[Discrete Logarithm, $\dlog$]
Let $G$ be a cyclic group with generator $g$. Given $h = g^x$, find $x$.
In $\Zp$ the security parameter is $n = \log p$ (equivalently $n = \log\abs{G}$).
\end{problem}

\begin{problem}[Computational Diffie--Hellman, $\cdh$]
Given $g^a$ and $g^b$ (but not $a,b$), compute $g^{ab}$.
\end{problem}

\begin{problem}[Decisional Diffie--Hellman, $\ddh$]
Distinguish the distribution $(g^a, g^b, g^{ab})$ from $(g^a, g^b, g^c)$ for independent
uniform $a,b,c$.
\end{problem}

\subsection{Lattice family}

A \emph{lattice} $\lat \subset \R^m$ is a discrete additive subgroup, i.e.\ the set of all
integer combinations of some linearly independent basis $\vB$. Write $\lamone(\lat)$ for the
length of a shortest nonzero vector and $\lambda_i$ for the successive minima.

\begin{problem}[Shortest / Shortest-Independent Vector, $\svp_\gamma$ / $\sivp_\gamma$]
Given a basis of $\lat$, find a nonzero $\vz\in\lat$ with $\norm{\vz}\le\gamma\cdot\lamone(\lat)$
(resp.\ $n$ independent vectors all of length $\le\gamma\cdot\lambda_n$). These are
\emph{worst-case} problems (over all lattices). Here $\gamma=\gamma(n)\ge 1$ is the
approximation factor; larger $\gamma$ makes the problem \emph{easier}.
\end{problem}

\begin{problem}[Short Integer Solution, $\sis_{n,m,q,\beta}$]
For uniformly random $\vA \in \Zq^{n\times m}$, find a nonzero $\vz \in \Z^m$ with
\[
   \vA\vz = \vzero \pmod q, \qquad \norm{\vz}\le\beta .
\]
Equivalently, find a short vector in the \emph{kernel lattice}
$\lat^{\perp}(\vA) = \{\vz\in\Z^m : \vA\vz = \vzero \bmod q\}$.
\end{problem}

\begin{problem}[Learning With Errors, $\lwe$]
For uniform $\vA$, secret $\vs\in\Zq^n$, and small error $\ve$, recover $\vs$ from
\[
   \vb = \vA^{t}\vs + \ve \pmod q .
\]
Geometrically this is \emph{bounded-distance decoding} on the LWE lattice
$\lat(\vA) = \{\vA^t\vs : \vs\in\Zq^n\} + q\Z^m = q\cdot\lat^\perp(\vA)^{*}$,
the dual of the SIS lattice up to the scale factor $q$.
\end{problem}

\begin{remark}[Search vs.\ decision]
Each problem has a \emph{search} version (find the secret) and a \emph{decision} version
(distinguish structured from uniform). For DLOG/factoring the decision version is degenerate,
so the interesting decision problems are $\ddh$ (for DH) and decision-$\lwe$ (for lattices).
\end{remark}

%==========================================================
\section{The Lattice Assumption and Its Structure}
%==========================================================

\begin{assumption}[Lattice assumption]
The complexity of solving lattice problems in $n$-dimensional lattices grows
superpolynomially (or exponentially) in $n$.
\end{assumption}

Three structural features make this assumption unusually robust. Features
\ref{sec:dim} and \ref{sec:rsr} are the ones absent from the DLOG world;
Feature~\ref{sec:wc-ac} is the headline property that motivates the whole field.

\subsection{Worst-case to average-case hardness}
\label{sec:wc-ac}

Cryptography needs \emph{average-case} hardness (random keys must be hard), but a
problem can be worst-case hard while being easy on average. The lattice family is special
because \emph{Ajtai's reduction} transports worst-case hardness into the average case:

\begin{theorem}[Ajtai; Regev; informal]
Average-case $\sis$ (resp.\ $\lwe$) is at least as hard as worst-case
$\sivp_\gamma$ / $\gapsvp_\gamma$ for some $\gamma=\poly(n)$:
\[
   \text{worst-case } \sivp_\gamma \reduces \text{average-case } \sis
   \qquad(\text{and similarly for } \lwe).
\]
\end{theorem}

\begin{keybox}[Why this matters]
The hardness of the \emph{cryptographic} object (SIS/LWE, defined over a random $\vA$) is
\emph{inherited} from a \emph{different}, older, independently studied problem (SIVP over
\emph{arbitrary} lattices). This is an \emph{external anchor}, not a self-loop --- which is
exactly what the DLOG family lacks.
\end{keybox}

\subsection{Hardness propagates across dimensions}
\label{sec:dim}

Lattice problems in dimension $n$ reduce to lattice problems in dimension $m>n$. The
reduction is a simple \emph{embedding}: given a basis $\vB$ of an $n$-dimensional lattice,
build the $m$-dimensional block lattice
\[
   \vB \;\Longrightarrow\;
   \begin{pmatrix} \vB & \mathbf 0 \\[2pt] \mathbf 0 & \mathbf D \end{pmatrix},
   \qquad \mathbf D = \text{large (``}\infty\text{'') diagonal block}.
\]
The new coordinates carry basis vectors so long that no short vector ever uses them;
hence the short vectors of the big lattice are exactly those of $\vB$ (padded with zeros,
which leaves the norm unchanged). Therefore a solver for dimension $m$ solves dimension $n$:
\[
   \underbrace{\text{dim-}n\text{ problem}}_{\text{smaller}}
   \reduces
   \underbrace{\text{dim-}m\text{ problem}}_{\text{larger},\; m>n}.
\]

\begin{keybox}[Two readings of one reduction]
\textbf{Solver direction:} a dimension-$m$ solver yields a dimension-$n$ solver
(high $\Rightarrow$ low).\\
\textbf{Hardness direction:} low-dimensional hardness implies high-dimensional hardness
(low $\Rightarrow$ high).
\[
  \text{dim-}m\ \text{\emph{solver}} \Rightarrow \text{dim-}n\ \text{\emph{solver}}
  \quad\Longleftrightarrow\quad
  \text{dim-}n\ \text{\emph{hard}} \Rightarrow \text{dim-}m\ \text{\emph{hard}}.
\]
\end{keybox}

The consequence is that the family indexed by dimension $n$ is \emph{connected}: hardness
in small, well-studied dimensions backstops hardness in all larger dimensions. One does not
assume each dimension independently; they form a single coherent chain.

\subsection{Random self-reducibility}
\label{sec:rsr}

Lattices (and the number-theoretic problems) are \emph{random self-reducible}: any fixed
instance can be re-randomized into a uniformly random one, so that solving random instances
solves every instance. Hence average-case and worst-case hardness coincide \emph{within the
problem}. (For SIS/LWE this is subsumed by the stronger external reduction of
\S\ref{sec:wc-ac}; the point here is that the re-randomization machinery exists and produces
\emph{genuinely uniform} instances --- the property that lets an average-case oracle be used
at all.)

%==========================================================
\section{Why DLOG Is Not Analogous}
%==========================================================

One may try to mimic the lattice assumption for DLOG.

\begin{assumption}[Alternative DLOG assumption]
$\dlog(p_n)$ is hard, where $p_n$ is the smallest prime $>2^n$. This is ``equivalent to a
worst-case family of problems indexed by $n$,'' but the definition is \emph{ad hoc}: the
choice ``smallest prime $>2^n$'' is arbitrary, reverse-engineered only to manufacture a
family. Lattices get their family \emph{for free} from the dimension.
\end{assumption}

\subsection{The Prime Number Theorem makes DLOG average-case}
\label{sec:pnt}

The natural way to state the DLOG assumption is average-case over a random prime.

\begin{theorem}[Prime Number Theorem]
$\pi(x) \sim \dfrac{x}{\ln x}$, where $\pi(x)$ counts primes $\le x$.
\end{theorem}

\begin{proposition}[Counting $n$-bit primes]
The number of $n$-bit primes is
\[
   \pi(2^n)-\pi(2^{n-1})
   \;\approx\; \frac{2^{n-1}}{n\ln 2}
   \;=\; \Theta\!\left(\frac{2^n}{n}\right),
\]
i.e.\ \emph{exponentially many}. Equivalently, a random $n$-bit integer is prime with
probability $\approx \dfrac{1}{n\ln 2} = \Theta(1/n)$.
\end{proposition}

\begin{keybox}[Exponentially many, polynomially dense]
\begin{itemize}[leftmargin=1.4em,nosep]
  \item \textbf{Exponential count} $\Theta(2^n/n)$ $\Rightarrow$ the key space is astronomically
        large; an adversary cannot enumerate primes.
  \item \textbf{Polynomial density} $\Theta(1/n)$ $\Rightarrow$ one can \emph{sample} a random
        $n$-bit prime in polynomial time ($\approx n$ primality tests).
\end{itemize}
This is why the DLOG assumption is stated as ``$\dlog$ is hard for \emph{random} $p$'':
there are exponentially many groups to choose from, $p$ is drawn uniformly among $n$-bit
primes, and the assumption is average-case over that choice.
\end{keybox}

\subsection{Hardness does not propagate across moduli}

This is the crucial disanalogy with \S\ref{sec:dim}. There is \emph{no known} reduction
between the DLOG problems of different primes:
\[
   \dlog(p_n) \;\overset{?}{\Longrightarrow}\; \dlog(p_{n+1}),
   \qquad\qquad
   \dlog(\Z_p^*) \;\overset{?}{\reduces}\; \dlog(\Z_q^*).
\]
The groups $\Z_{p}^*$ and $\Z_{q}^*$ have unrelated orders and structure; there is no
padding-style embedding of one into the other. The primes are \emph{isolated islands}: each
$\dlog(p_n)$ is a \emph{separate} assumption, and random self-reducibility
(RSR) acts only \emph{within a fixed group} --- it certifies average-case hardness for a
given $p$ but \textbf{provides no guidance on how to choose $p$}, nor any bridge between
different $p$. The ``family'' has members that do not talk to each other.

\begin{remark}[Intuition is not a reduction]
It \emph{feels} intuitive that larger primes should be at least as hard. But the ``$?$''
above is literal: we have no construction turning one into the other, unlike the explicit,
provable lattice embedding. Cryptography wants theorems, not vibes.
\end{remark}

\subsection{Further fragilities of DLOG hardness}

\paragraph{Representation dependence.}
The group $(\Z_p^*,\cdot)$ is isomorphic to $(\Z_{p-1},+)$, yet ``DLOG'' in the additive
representation is \emph{trivial} (it is division modulo $p-1$). The \emph{same} group under a
different natural encoding flips DLOG from hard to easy. Hardness is thus a property of the
\emph{representation}, not of the group --- whereas lattice hardness is an intrinsic
geometric property, immune to relabeling.

\paragraph{Group dependence.}
Different natural groups ($\Z_p^*$, $\Z_{pq}^*$, elliptic curves, extension fields) exhibit
\emph{widely} different hardness, some with subexponential attacks and some without. There is
no single ``DLOG is hard''; there is a scattered zoo of group-specific assumptions.

\paragraph{One-directional CRT reduction.}
The Chinese Remainder Theorem gives $\Z_{pq}\cong\Z_p\times\Z_q$, hence
\[
   \dlog(\Z_p^*) \;\Longrightarrow\; \dlog(\Z_{pq}^*)
\]
(a composite-group solver projects down to a prime component, so the composite is at least
as hard as the prime). But the \emph{reverse} direction requires knowing $p,q$ separately,
i.e.\ \emph{factoring} $pq$. The two directions are asymmetric, and one smuggles in a second
assumption.

%==========================================================
\section{Search vs.\ Decision Reductions}
%==========================================================

Both families have search and decision versions. The contrast lies in whether the two are
\emph{equivalent}.

\subsection{LWE: search \texorpdfstring{$\Longleftrightarrow$}{<=>} decision}

\paragraph{Easy direction (decision \texorpdfstring{$\reduces$}{<=} search).}
A search oracle trivially gives a decider. Given $(\vA,\vb)$, run search to get a candidate
$\vs'$ and test the residual:
\[
   \norm{\vb - \vA^t\vs' \bmod q}\ \text{small} \;\Rightarrow\; \text{real LWE},
   \qquad
   \text{large} \;\Rightarrow\; \text{uniform}.
\]
Checking $\vs'$ costs one matrix multiply; this verifiability is what makes the direction
trivial. (Note: a small error bound guarantees the answer is \emph{unique and verifiable},
not that it is \emph{easy to find} --- recovering $\vs$ is still hard because Gaussian
elimination amplifies the error.)

\paragraph{Hard direction (search \texorpdfstring{$\reduces$}{<=} decision).}
Here a decision oracle $\mathcal D$ (one bit: ``real vs.\ uniform'') is bootstrapped into
full recovery of $\vs$, one coordinate at a time. To learn coordinate $s_1$:
\begin{enumerate}[leftmargin=1.6em,nosep]
  \item \textbf{Guess} a value $g\in\Zq$ for $s_1$.
  \item \textbf{Modify} the instance using a \emph{fresh uniform} $r\in\Zq$, so that the
        residual contribution of the first coordinate becomes $(s_1-g)\cdot r$.
  \item \textbf{Query} $\mathcal D$ on the modified instance.
\end{enumerate}
If the guess is correct, $s_1-g=0$ and the instance stays valid LWE ($\mathcal D$:
``real''). If the guess is wrong, $s_1-g\neq 0$; since \emph{nonzero $\times$ uniform
$=$ uniform} in the field $\Zq$, the term $(s_1-g)r$ is uniform and scrambles the sample to
look random ($\mathcal D$: ``uniform''). Thus
\[
   \mathcal D(\text{modified instance}) = \text{``real''} \iff g = s_1 .
\]
Trying all $q$ values per coordinate over all $n$ coordinates uses $nq=\poly(n)$ queries
(for $q=\poly(n)$), recovering $\vs$ entirely.

\begin{keybox}[Why the randomizer $r$ is essential]
Without $r$, subtracting the guess merely \emph{shifts the secret}, and shifted-secret LWE is
\emph{still valid LWE}; the oracle answers ``real'' for right \emph{and} wrong guesses,
leaking nothing. The random $r$ converts a wrong guess into \emph{uniform} --- the only
signal a decision oracle can recognize. (Same re-randomization principle as in the
worst-case/average-case reductions: feed the oracle instances from exactly the distribution
it recognizes.)
\end{keybox}

\begin{theorem}[Search--decision equivalence for LWE]
For appropriate parameters, search-$\lwe$ and decision-$\lwe$ are polynomially equivalent.
\end{theorem}

\begin{keybox}[Consequence]
One-wayness (search) and pseudorandomness (decision) rest on the \emph{same} assumption.
Pseudorandomness comes \emph{for free}, which is what lets LWE samples mask messages and
yield semantically secure encryption (a PRG-level primitive, well above a bare one-way
function).
\end{keybox}

\subsection{Diffie--Hellman: search and decision are \emph{not} equivalent}

\paragraph{Easy direction (decision \texorpdfstring{$\reduces$}{<=} search, i.e.\ $\ddh\reduces\cdh$).}
A $\cdh$ solver gives a $\ddh$ decider: given $(g^a,g^b,Z)$, compute $g^{ab}$ and test
$Z\overset{?}{=}g^{ab}$. So $\ddh$ is no harder than $\cdh$.

\paragraph{The reverse ($\ddh\Longrightarrow\cdh$) is not available in general.}
Unlike LWE, one \emph{cannot} generally turn a $\ddh$ distinguisher into a $\cdh$ solver.
In fact the two problems are \emph{provably inequivalent} in important groups:

\begin{keybox}[Pairing groups: DDH easy, CDH hard]
On \emph{pairing-friendly elliptic curves} there is a bilinear map $e$ with
$e(g^a,g^b)=e(g,g)^{ab}$. This lets one \emph{decide} DDH efficiently
(compare $e(g^a,g^b)$ with $e(g,Z)$), yet \emph{computing} $g^{ab}$ remains hard.
Hence \textbf{DDH is broken while CDH survives} --- the search/decision gap is real,
not a missing proof.
\end{keybox}

\paragraph{Pairing-based cryptography.}
This very gap founded an entire subfield: schemes that \emph{require} DDH to be easy (for
functionality via the pairing) while relying on CDH-type problems (or the Bilinear
Diffie--Hellman problem) for security. Landmark examples include one-round three-party key
agreement (Joux), identity-based encryption (Boneh--Franklin), and short signatures
(Boneh--Lynn--Shacham). None of this would be possible if DDH and CDH were equivalent.

\subsection{The DH hardness ladder}

The discrete-log problems form a ladder in which the ``easy'' (downward) direction always
holds, while the ``hard'' (upward) direction is either open or false:
\[
   \ddh \reduces \cdh \reduces \dlog .
\]
\begin{itemize}[leftmargin=1.5em,nosep]
  \item $\dlog\Rightarrow\cdh$ (recover the exponent, then exponentiate) --- easy.
  \item $\cdh\Rightarrow\ddh$ (compute and compare) --- easy.
  \item $\cdh\Rightarrow\dlog$? \textbf{Open in general} (partial results: Maurer--Wolf,
        under auxiliary conditions).
  \item $\ddh\Rightarrow\cdh$? \textbf{False in general} (pairing groups).
\end{itemize}
Analogously, whether breaking \emph{RSA} is as hard as \emph{factoring} is also open.

%==========================================================
\section{Summary}
%==========================================================

\begin{table}[h]
\centering
\renewcommand{\arraystretch}{1.35}
\begin{tabular}{@{}p{0.34\textwidth} p{0.30\textwidth} p{0.30\textwidth}@{}}
\toprule
\textbf{Structural feature} & \textbf{DLOG / Diffie--Hellman} & \textbf{Lattices / LWE} \\
\midrule
Average-case over which object?
  & random prime $p$ (via PNT: $\Theta(2^n/n)$ primes)
  & random matrix $\vA$ \\
Hardness across the family
  & \textbf{isolated}: $\dlog(p_n)\overset{?}{\Rightarrow}\dlog(p_{n+1})$ unknown
  & \textbf{connected}: dim-$n$ $\reduces$ dim-$m$ by embedding \\
Random self-reducibility
  & yes, but only \emph{within a fixed group}; no guidance on choosing $p$
  & yes, producing genuinely uniform instances \\
External worst-case anchor
  & none (self-loop only)
  & yes (Ajtai: worst-case SIVP $\reduces$ SIS/LWE) \\
Representation dependence
  & \textbf{yes} ($\Z_p^*\!\cong\!\Z_{p-1}$: additive DLOG trivial)
  & no (intrinsic geometry) \\
Search $\reduces$ decision (easy)
  & yes ($\ddh\reduces\cdh$)
  & yes (check the residual) \\
Decision $\reduces$ search (hard)
  & \textbf{false in general} (pairings: DDH easy, CDH hard)
  & \textbf{proven} (guess-and-test with randomizer $r$) \\
Search $\Leftrightarrow$ decision
  & \textbf{no} (separate, strictly stronger, sometimes broken)
  & \textbf{yes} (equivalent) \\
Quantum status
  & broken by Shor
  & no known quantum speedup \\
\bottomrule
\end{tabular}
\caption{Reduction structure of the two assumption families.}
\end{table}

\begin{keybox}[Takeaways]
\begin{enumerate}[leftmargin=1.5em,nosep]
  \item Lattice hardness is a \emph{connected, intrinsic, externally anchored} family;
        hardness propagates from low to high dimension by an explicit embedding.
  \item DLOG hardness is a \emph{disconnected, representation-dependent} patchwork; there is
        no reduction across moduli, and RSR gives no guidance on choosing the group.
  \item The average-case nature of DLOG rests on the PNT: exponentially many primes
        (huge key space) that are polynomially dense (efficient sampling).
  \item For LWE, \emph{search $\Leftrightarrow$ decision} is a theorem, so pseudorandomness
        is free. For Diffie--Hellman the directions split: $\ddh\reduces\cdh$ is easy, but
        $\ddh\Rightarrow\cdh$ is false in pairing groups --- a gap that \emph{founded}
        pairing-based cryptography.
\end{enumerate}
\end{keybox}

\end{document}
